%%%%%%%%%%%%%%%%%%%%%%%%%%%%%%%%%%%%%%%%%%%%%%%%%%%%%%%%%%%%%%%%%%%%%%%%%%%%
% 1. FONTS: Times-style text and math font.
%    Alternatives (swap this block):
%      Libertine:  \usepackage{libertine} + \usepackage[libertine]{newtxmath}
%      XCharter:   \usepackage{xcharter}
%%%%%%%%%%%%%%%%%%%%%%%%%%%%%%%%%%%%%%%%%%%%%%%%%%%%%%%%%%%%%%%%%%%%%%%%%%%%
\usepackage{newtxtext}

%%%%%%%%%%%%%%%%%%%%%%%%%%%%%%%%%%%%%%%%%%%%%%%%%%%%%%%%%%%%%%%%%%%%%%%%%%%%
% 2. MICROTYPOGRAPHY: subtle spacing tweaks for cleaner justified text.
%%%%%%%%%%%%%%%%%%%%%%%%%%%%%%%%%%%%%%%%%%%%%%%%%%%%%%%%%%%%%%%%%%%%%%%%%%%%
\usepackage{microtype}

%%%%%%%%%%%%%%%%%%%%%%%%%%%%%%%%%%%%%%%%%%%%%%%%%%%%%%%%%%%%%%%%%%%%%%%%%%%%
% 3. LANGUAGE & ENCODING: special characters + hyphenation rules.
%    Swap "english" for your language if needed, e.g.:
%      \usepackage[ngerman]{babel}   % German
%      \usepackage[french]{babel}    % French
%    This automatically updates hyphenation patterns and auto-generated
%    text like "Figure"/"Abbildung" without touching anything else here.
%%%%%%%%%%%%%%%%%%%%%%%%%%%%%%%%%%%%%%%%%%%%%%%%%%%%%%%%%%%%%%%%%%%%%%%%%%%%
\usepackage[utf8]{inputenc}
\usepackage[T1]{fontenc}
\usepackage{lmodern}
\usepackage[english]{babel}
\hyphenpenalty=1000
\exhyphenpenalty=50
\tolerance=1000

%%%%%%%%%%%%%%%%%%%%%%%%%%%%%%%%%%%%%%%%%%%%%%%%%%%%%%%%%%%%%%%%%%%%%%%%%%%%
% 4. MATH: equations, symbols, theorem environments, bold vectors.
%%%%%%%%%%%%%%%%%%%%%%%%%%%%%%%%%%%%%%%%%%%%%%%%%%%%%%%%%%%%%%%%%%%%%%%%%%%%
\usepackage{amsmath}
\usepackage{amssymb}
\usepackage{amsthm}
\usepackage{mathtools}
\usepackage{bm}
\usepackage{newtxmath} % load after amsmath; swap if using an alt font (see above)
\let\openbox\relax     % avoids a naming clash between amsthm and another package

%%%%%%%%%%%%%%%%%%%%%%%%%%%%%%%%%%%%%%%%%%%%%%%%%%%%%%%%%%%%%%%%%%%%%%%%%%%%
% 5. GRAPHICS, FIGURES & TABLES
%
%    TIP: Building tables by hand is tedious and error-prone, especially
%    regression output tables (coefficients, standard errors, significance
%    stars, footnotes). Recommended workflow instead:
%      1. Run your regression in R (or Python/Stata) and export the raw
%         output, or just take a screenshot of the R console/table.
%      2. Paste the R output or screenshot into an AI assistant and ask it
%         to generate a LaTeX table using booktabs (\toprule, \midrule,
%         \bottomrule) that matches your regression results.
%      3. Compile the generated code as-is, then manually tweak only small
%         details afterward: column widths, number formatting, footnote
%         text for significance levels (e.g. * p<0.10, ** p<0.05), or
%         table width via adjustbox if it overflows the page margin.
%    This is far faster than manually transcribing coefficients and gets
%    you a clean, publication-style table (similar to R's stargazer or
%    modelsummary output) without writing raw tabular code by hand.
%%%%%%%%%%%%%%%%%%%%%%%%%%%%%%%%%%%%%%%%%%%%%%%%%%%%%%%%%%%%%%%%%%%%%%%%%%%%
\usepackage{graphicx}
\usepackage{float}
\usepackage{caption}
\usepackage{subcaption}
\usepackage{booktabs}
\usepackage{array}
\usepackage{multirow}
\usepackage{rotating}
\usepackage{lscape}
\usepackage{pdflscape}
\usepackage{adjustbox}
\usepackage{tcolorbox}

%%%%%%%%%%%%%%%%%%%%%%%%%%%%%%%%%%%%%%%%%%%%%%%%%%%%%%%%%%%%%%%%%%%%%%%%%%%%
% 6. UNITS (OPTIONAL): only needed for physical units or numeric tables
%    that require consistent decimal alignment, e.g. "3.14 kg".
%    Comment out if you don't work with measurements/units.
%%%%%%%%%%%%%%%%%%%%%%%%%%%%%%%%%%%%%%%%%%%%%%%%%%%%%%%%%%%%%%%%%%%%%%%%%%%%
% \usepackage{siunitx}
% \sisetup{detect-all}

%%%%%%%%%%%%%%%%%%%%%%%%%%%%%%%%%%%%%%%%%%%%%%%%%%%%%%%%%%%%%%%%%%%%%%%%%%%%
% 7. TEXT DECORATION & LISTS
%%%%%%%%%%%%%%%%%%%%%%%%%%%%%%%%%%%%%%%%%%%%%%%%%%%%%%%%%%%%%%%%%%%%%%%%%%%%
\usepackage{enumitem}
\usepackage[normalem]{ulem} % enables \sout{} strikethrough

%%%%%%%%%%%%%%%%%%%%%%%%%%%%%%%%%%%%%%%%%%%%%%%%%%%%%%%%%%%%%%%%%%%%%%%%%%%%
% 8. DIAGRAMS (OPTIONAL): TikZ for custom diagrams/flowcharts drawn in code.
%
%    TIP: Writing TikZ by hand is painful, especially node positioning and
%    arrows. Recommended workflow instead:
%      1. Describe the diagram to an AI assistant (e.g. "causal diagram
%         showing X affects Y through mediator Z, with labeled arrows")
%         and ask it to generate the TikZ code.
%      2. Compile the generated code as-is to see the rough layout.
%      3. Manually tweak only the fine details afterward: node coordinates
%         (via the "positioning" library, e.g. right=of, below=of),
%         arrow styles, and spacing values, since these are the parts
%         AI often gets slightly off but are easy to nudge by eye.
%    TikZ is very flexible this way: flowcharts, causal mechanisms,
%    timelines, or statistical illustrations all render as crisp vector
%    graphics that scale perfectly in the final PDF, unlike a screenshot.
%    Uncomment only if you actually plan to draw diagrams in LaTeX.
%%%%%%%%%%%%%%%%%%%%%%%%%%%%%%%%%%%%%%%%%%%%%%%%%%%%%%%%%%%%%%%%%%%%%%%%%%%%
\usepackage{tikz}
\usetikzlibrary{arrows.meta,positioning,fit,calc,decorations.pathreplacing}

%%%%%%%%%%%%%%%%%%%%%%%%%%%%%%%%%%%%%%%%%%%%%%%%%%%%%%%%%%%%%%%%%%%%%%%%%%%%
% 9. COLORS: custom color definitions used in tables/highlights.
%%%%%%%%%%%%%%%%%%%%%%%%%%%%%%%%%%%%%%%%%%%%%%%%%%%%%%%%%%%%%%%%%%%%%%%%%%%%
\usepackage{xcolor}
\usepackage{colortbl}
\definecolor{regblue}{RGB}{220, 230, 242}
\definecolor{headerblue}{RGB}{169, 196, 223}
\definecolor{emptycell}{RGB}{245, 245, 245}

%%%%%%%%%%%%%%%%%%%%%%%%%%%%%%%%%%%%%%%%%%%%%%%%%%%%%%%%%%%%%%%%%%%%%%%%%%%%
% 10. BIBLIOGRAPHY: author-year citations (natbib/plainnat = Harvard style,
%     common in economics).
%       See folder Front Matter and then document "bibliography_setup.tex" for      explaination how it works.
%       If you want to edit, edit in "bibliography_setup.tex".
%%%%%%%%%%%%%%%%%%%%%%%%%%%%%%%%%%%%%%%%%%%%%%%%%%%%%%%%%%%%%%%%%%%%%%%%%%%%
%%%%%%%%%%%%%%%%%%%%%%%%%%%%%%%%%%%%%%%%%%%%%%%%%%%%%%%%%%%%%%%%%%%%%%%%%%%%
% BIBLIOGRAPHY: author-year citations (natbib/plainnat = Harvard style,
%     common in economics).
%
%     HOW CITATIONS WORK (mechanism explained):
%     - All sources live in a separate plain-text file, e.g. references.bib,
%       NOT in this .tex file. Each entry has a unique key (e.g. smith2020)
%       plus fields like author, title, year, journal.
%     - \bibliography{references} at the END of your document (without the
%       .bib extension) tells LaTeX which file to pull from.
%     - \bibliographystyle{plainnat} below controls the FORMAT of both the
%       in-text citations and the final reference list. You never type the
%       full reference yourself; it's auto-generated from the .bib fields.
%     - \citep{smith2020}  ->  "(Smith, 2020)"   use when the name is NOT
%       part of your sentence grammar.
%     - \citet{smith2020}  ->  "Smith (2020)"    use when the name flows
%       into the sentence, e.g. "As \citet{smith2020} argues..."
%
%     ABBREVIATING A FREQUENTLY-CITED SOURCE:
%     If you cite the same paper repeatedly (e.g. a key theory paper),
%     define a shortcut once in the preamble:
%       \newcommand{\smithcite}{\citep{smith2020}}
%     Then just type \smithcite throughout instead of retyping the full
%     \citep{smith2020} each time. Useful if you later want to change how
%     that one source is cited everywhere at once.
%
%     COSMETIC APA-LIKE FALLBACK (stay on natbib/bibtex, no engine switch):
%     If a professor asks for "APA-ish" formatting but you don't want to
%     deal with biblatex+biber, just change the style below to:
%       \bibliographystyle{apalike}
%     Everything else (citep/citet, bibtex compilation) stays identical.
%
%     FULL APA 7 / CHICAGO AUTHOR-DATE (more accurate, bigger switch):
%     Requires replacing natbib entirely with biblatex + biber backend and
%     renaming citation commands (citep -> autocite/parencite,
%     citet -> textcite). Only worth it if the professor requires the
%     precise official style rather than something "close enough".
%%%%%%%%%%%%%%%%%%%%%%%%%%%%%%%%%%%%%%%%%%%%%%%%%%%%%%%%%%%%%%%%%%%%%%%%%%%%
\usepackage[round,authoryear]{natbib}
\bibliographystyle{plainnat}  % change to apalike for a quick APA-like look
\setlength{\bibhang}{1em}
\renewcommand{\bibfont}{\small}
\setlength{\bibsep}{1.2em}

%%%%%%%%%%%%%%%%%%%%%%%%%%%%%%%%%%%%%%%%%%%%%%%%%%%%%%%%%%%%%%%%%%%%%%%%%%%%
% 11. PAGE LAYOUT
%%%%%%%%%%%%%%%%%%%%%%%%%%%%%%%%%%%%%%%%%%%%%%%%%%%%%%%%%%%%%%%%%%%%%%%%%%%%
\usepackage[a4paper,left=2.5cm,right=2.5cm,top=2.5cm,bottom=2.5cm]{geometry}
\usepackage{setspace}
\onehalfspacing
\setlength{\parindent}{0pt}
\setlength{\parskip}{1em}
\usepackage{multicol}

%%%%%%%%%%%%%%%%%%%%%%%%%%%%%%%%%%%%%%%%%%%%%%%%%%%%%%%%%%%%%%%%%%%%%%%%%%%%
% 12. FOOTNOTES: smaller font, single line spacing.
%%%%%%%%%%%%%%%%%%%%%%%%%%%%%%%%%%%%%%%%%%%%%%%%%%%%%%%%%%%%%%%%%%%%%%%%%%%%
\usepackage[hang, flushmargin]{footmisc}
\renewcommand{\footnotelayout}{\setstretch{1.0}\footnotesize}

%%%%%%%%%%%%%%%%%%%%%%%%%%%%%%%%%%%%%%%%%%%%%%%%%%%%%%%%%%%%%%%%%%%%%%%%%%%%
% 13. HEADER & FOOTER
%%%%%%%%%%%%%%%%%%%%%%%%%%%%%%%%%%%%%%%%%%%%%%%%%%%%%%%%%%%%%%%%%%%%%%%%%%%%
\usepackage{fancyhdr}
\setlength{\headheight}{14pt}
\pagestyle{fancy}
\fancyhf{}
\fancyhead[L]{[Insert Type of Document]} %Adjust
\fancyhead[R]{\today}
\fancyfoot[C]{\thepage}

%%%%%%%%%%%%%%%%%%%%%%%%%%%%%%%%%%%%%%%%%%%%%%%%%%%%%%%%%%%%%%%%%%%%%%%%%%%%
% 14. CLICKABLE LINKS & SMART REFERENCES
%       cleveref enables \cref{} for automatic "Figure 3" style references
%       instead of manually typing "Figure~\ref{}".
%       Adjust for a different color (e.g. blue), research options before changing.
%%%%%%%%%%%%%%%%%%%%%%%%%%%%%%%%%%%%%%%%%%%%%%%%%%%%%%%%%%%%%%%%%%%%%%%%%%%%
\usepackage[
    colorlinks=true,
    linkcolor=black,                    
    citecolor=black,
    urlcolor=blue,
    pdftitle={Your Document},
    pdfpagemode=FullScreen,
]{hyperref}
\usepackage{cleveref}
\usepackage{xurl}

%%%%%%%%%%%%%%%%%%%%%%%%%%%%%%%%%%%%%%%%%%%%%%%%%%%%%%%%%%%%%%%%%%%%%%%%%%%%
% 15. CUSTOM SHORTCUTS: your own math notation commands, logo insert or what else.
%       Add all your other \newcommand's below here so its tidy.
%       Example: \E{X} produces E[X] (expected value).
%%%%%%%%%%%%%%%%%%%%%%%%%%%%%%%%%%%%%%%%%%%%%%%%%%%%%%%%%%%%%%%%%%%%%%%%%%%%
\newcommand{\UniversityLogo}{images/university_logo.png}      % path to logo file
\newcommand{\UniversityLogoScale}{0.015}                      % scale of the logo
\newcommand{\E}[1]{\mathbb{E}\left[#1\right]}
\newcommand{\Var}[1]{\mathrm{Var}\left[#1\right]}
\newcommand{\B}{\beta}

